% CVPR 2026 author-kit document setup; camera-ready mode displays all authors.
\documentclass[10pt,twocolumn,letterpaper]{article}
\usepackage{cvpr}
%% This file contains a number of tweaks that are typically applied to the main document.
%% They are not enabled by default, but can be enabled by uncommenting the relevant lines.


%%
%% Official support for 2-column teaser figure
%%
\usepackage{cuted} %< for strip environment

%%
%% For auto-referencing of figures (see fig/teaser.tex and \label{} therein) 
%%
\usepackage{currfile} %< for \currfilebase macro

%%
%% caption and (global) spacing overrides
%%
\usepackage{caption} %< for "\captionof" commands (teaser)
% \setlength{\abovecaptionskip}{.5em} % space above caption
% \setlength{\belowcaptionskip}{1em} % space below caption

%%
%% Inline annotations; for predefined colors, refer to "dvipsnames" in the xcolor package:
%% https://tinyurl.com/overleaf-colors
%%
\newcommand{\red}[1]{{\color{red}#1}}
\newcommand{\todo}[1]{{\color{red}#1}}
\newcommand{\TODO}[1]{\textbf{\color{red}[TODO: #1]}}
%%
%% Disable for camera ready / submission by uncommenting these lines:
%%
% \renewcommand{\TODO}[1]{}
% \renewcommand{\todo}[1]{#1}

%%
%% Work harder in optimizing text layout. Typically shrinks text by 1/6 of page, enable
%% it at the very end of the writing process, when you are just above the page limit.
%%
% \usepackage{microtype}

%%
%% Fine-tune paragraph spacing.
%%
% \renewcommand{\paragraph}[1]{\vspace{.5em}\noindent\textbf{#1.}}

%%
%% Globally adjusts space between figure and caption.
%%
% \setlength{\abovecaptionskip}{.5em}

%%
%% Allows "the use of \paper to refer to the project name"
%% with automatic management of space at the end of the word.
%%
% \usepackage{xspace}
% \newcommand{\paper}{ProjectName\xspace}

%%
%% Commonly used math definitions.
%%
% \DeclareMathOperator*{\argmin}{arg\,min}
% \DeclareMathOperator*{\argmax}{arg\,max}

%%
%% Tighten underline.
%%
% \usepackage{soul}
% \setuldepth{foobar}

%%
%% Reconfigure paragraph (tweaks spacing)
%%

\definecolor{cvprblue}{rgb}{0.21,0.49,0.74}
\usepackage[pagebackref,breaklinks,colorlinks,allcolors=cvprblue]{hyperref}
\def\paperID{}
\def\confName{CVPR}
\def\confYear{2026}
\input{generated/active_study/settings.tex}
\title{Teacher-Calibrated Perceptual Diversity for One-Step Diffusion}
\author{Le Duc Dong \quad Nguyen Huu Thuat \quad Pham Viet Hoang \quad La Nguyen Gia Hy\\CE7454 Project 2}
\begin{document}
\maketitle
\begin{abstract}
We investigate direct one-step text-to-image generation with a fixed rank-32
Stable Diffusion student and a 36,000-step training budget. Our DMD-based
proposal combines distribution matching, perceptual regression and matching
the teacher's variation across noise seeds for the same caption. Evaluation
uses all 5,000 COCO-2017 validation images and matched captions and noise.
The completed proposal achieves FID \ActiveFID{} and CLIP \ActiveCLIP{} at one denoiser
evaluation, compared with FID \BaselineDMDFID{} and CLIP \BaselineDMDCLIP{} for our DMD
baseline. \ActiveAblationNote{} We provide measured baseline
results, selected qualitative examples and explicit experimental limitations.
\end{abstract}

\section{Introduction}
Diffusion models generate detailed images but require repeated denoiser
evaluations. Distillation transfers a pretrained teacher into a generator
with a smaller inference budget. Our goal is a practical one-step student
that preserves visual detail and useful variation under a limited adapter
budget. A low reconstruction loss alone is insufficient: blurry or repetitive
images can fit a restrictive target without capturing the teacher distribution.

We study whether perceptual regression and teacher-calibrated within-prompt
variation help a DMD-based LoRA student. The proposal matches the perceptual
distance between two student images to the distance between their paired
teacher outputs. This target avoids rewarding image separation without a
teacher reference. The contribution is a controlled image-level extension
under a fixed backbone, student capacity and total training-step budget.

\begin{figure*}[t]
\centering
\includegraphics[width=\textwidth]{generated/active_study/selected_samples.pdf}
\caption{Selected qualitative examples from completed 36,000-step students.
Each row uses the identical recorded caption and initial evaluation noise
across methods. All students use one step / one NFE; the teacher is a
64-step / 128-NFE reference. SwiftBrush uses teacher CFG 4.5 and FP32;
other students use CFG 7.5 and BF16. The selected subjects are a dog in flowers,
teddy bears and a lighthouse. Dog/lighthouse use fixed preset prompts;
the teddy-bear row is a validation example. Full-set metrics
use all 5,000 validation outputs. Baseline names denote our LoRA adaptations.}
\label{fig:demo}
\end{figure*}

\section{Related work}
Latent diffusion enables high-resolution generation using a pretrained VAE
and a denoiser operating on compressed image latents~\cite{rombach2022ldm}.
LoRA reduces the number of trained parameters~\cite{hu2022lora}. Progressive
distillation repeatedly replaces two teacher transitions with one student
transition~\cite{salimans2022progressive}. Consistency models learn predictions
that agree along a probability-flow trajectory~\cite{song2023consistency}.
InstaFlow uses reflow to straighten noise--image couplings before one-step
distillation~\cite{liu2024instaflow}. These objectives provide complementary
baselines for a common one-step inference budget.

DMD combines distribution matching with regression to teacher-generated
outputs~\cite{yin2024dmd}; SwiftBrush trains a generator using variational
score distillation~\cite{nguyen2024swiftbrush}. Perceptual regression is
already part of original DMD. Relational diffusion distillation also studies
cross-sample information~\cite{feng2024rdd}. We therefore evaluate a specific
same-caption teacher-distance objective in our restricted setting, without
claiming the first use of perceptual or relational distillation. Our baseline
implementations use reduced-capacity adapters and differ from the published
full-model recipes; their measurements are local experimental results.

\section{Method}
\paragraph{One-step generator and distribution matching.}
The student $G_\theta(z,c)$ maps Gaussian noise $z$ and a caption $c$ to a
clean VAE latent in one conditional U-Net call. Teacher and fake-score
networks supply clean estimates $\hat x_0^T$ and $\hat x_0^F$ from a noisy
student latent. The DMD gradient direction is
\begin{equation}
g=\frac{\hat x_0^F-\hat x_0^T}
 {\operatorname{mean}(|G_\theta-\hat x_0^T|)+\epsilon}.
\end{equation}
The implementation injects the stopped direction into the generator and
fits the auxiliary scorer with epsilon denoising on detached student outputs.
Both generator and scorer use rank-32 adapters. The VAE and text encoder
remain frozen. Guidance is incorporated during training; student inference
uses a single conditional branch with no additional CFG pass.

\paragraph{Teacher-calibrated perceptual objective.}
Let $I_i=D(G_\theta(z_i,c))$ and $J_i=D(T(z_i,c))$, where $D$ is the VAE
decoder and $T$ denotes a cached teacher trajectory. RGB values are clamped
to $[-1,1]$ for frozen LPIPS-VGG distance $d$~\cite{zhang2018lpips}. We use
\begin{align}
 L_{\mathrm{reg}}&=\tfrac12\sum_{i=1}^{2}d(I_i,J_i),\\
 L_{\mathrm{div}}&=\big[d(I_1,I_2)-\operatorname{sg}(d(J_1,J_2))\big]^2,\\
 L&=L_{\mathrm{DMD}}+0.25L_{\mathrm{reg}}+0.1L_{\mathrm{div}}.
\end{align}
The two noise seeds are independent, while the caption is shared. Both
student branches receive gradients. Teacher outputs and LPIPS weights
are frozen; decoder backpropagation supplies gradients to the student latent.
Only the student denoiser and decoder are needed at inference.

\paragraph{Controlled ablations.}
Four variants cross latent-MSE versus LPIPS regression with diversity weight
$\lambda_d\in\{0,0.1\}$. They share the exact same DMD initialization,
teacher-pair bank, captions, optimizer state, training seed and effective
batch size. Each follows a 36,000-step pipeline: the common DMD initialization
phase accounts for half the budget, followed by the loss-specific phase.
The grouped latent-MSE control is required to separate changes in the
anchor protocol from the effects of the proposed losses.

\section{Experimental protocol}
\paragraph{Data and evaluation coverage.}
We use SD1.5 with $512\times512$ outputs. COCO-2017 captions are ordered
by image ID, selecting the lowest-ID caption per image~\cite{chen2015coco}.
Training uses disjoint pools of 50,000 and 68,287 training-image captions.
Together these cover the 118,287 available training images. The proposed
DMD-family objectives use captions, random noise and generated teacher
anchors, with no ground-truth image supervision in those losses.
All 5,000 official validation images are held out from student training.
Reference images use bicubic resize and center crop. Evaluation covers every
official validation image with its selected caption.

Every final quality evaluation generates one image for each of the 5,000
validation captions, using the same evaluation noise stream (seed 12345).
FID and Inception Score (IS) measure distribution-level generation quality.
CLIP uses ViT-B/32 and reports mean $100\max(\cos(I,c),0)$. These are
full COCO-2017 validation results under our caption/preprocessing protocol,
which differs from published COCO-2014 benchmark protocols.

\paragraph{Capacity, updates and compute.}
All students have rank-32 LoRA generators with 6.38 million trainable
parameters, effective batch 16, training seed 42, and 36,000
generator updates. All student comparisons use one step / one NFE.
Our DMD study uses microbatch two with eight accumulation passes, AdamW,
generator/scorer learning rates $10^{-4}$, 100 warmup steps, EMA 0.995 and
one scorer update per generator update. Its shared bank has 5,000 teacher
images from 2,500 prompts with two noise seeds each. The frozen teacher uses
64 DDIM steps and CFG 7.5, costing 640,000 denoiser evaluations to create
the shared bank once. Decoder and LPIPS gradients add training compute;
equal generator updates do not imply equal total compute.

PD uses the epsilon-backbone progressive schedule; consistency uses a VP
boundary blend and latent L2. DMD uses normalized score matching with latent
regression, while InstaFlow uses the VP-to-velocity adaptation and cached
reflow couplings. SwiftBrush retains a rank-32 student with its recorded
optimizer/scorer settings, teacher CFG 4.5 and FP32 precision. Other students
use teacher CFG 7.5 and BF16. All methods appear in one comparison, with
these recorded settings disclosed. Concurrent GPU workloads prevent controlled latency claims; NFE
provides the inference-budget comparison.

\paragraph{Diversity and visual selection.}
Diversity is measured separately on 128 held-out captions with four fixed
seeds each, averaging the six within-caption LPIPS-Alex distances. AlexNet
differs from the VGG training metric. This 512-image diversity test is distinct
from the full 5,000-image quality evaluation. Figure~\ref{fig:demo} uses
three selected subjects with identical row captions and noise across methods.
Two use preset prompts; the teddy-bear example comes from validation.

\section{Results and analysis}
\begin{table}[t]
\centering\scriptsize
\resizebox{\linewidth}{!}{\input{generated/active_study/quality.tex}}
\caption{Full-validation quality: 5,000 generated images and 5,000 real
references. Student rows use the full 36,000-step budget. Bold marks the best
student value within this comparison. SwiftBrush uses teacher CFG 4.5/FP32 and other students
use CFG 7.5/BF16; the teacher has a longer trajectory. All baseline names refer to
our implementations.}
\label{tab:quality}
\end{table}
\input{generated/active_study/measured.tex}

\begin{table}[t]
\centering\scriptsize
\resizebox{\linewidth}{!}{\input{generated/active_study/ablation.tex}}
\caption{Factorial study with a shared initialization and 36,000 total
training steps. Diversity is pairwise LPIPS-Alex distance across four seeds
for each of 128 prompts. Bold marks the lowest ablation FID; a larger
perceptual distance does not itself imply higher quality.}
\label{tab:ablation}
\end{table}
\input{generated/active_study/prompt_control.tex}
\input{generated/active_study/interpolation_figure.tex}

\paragraph{Visual findings and failure modes.}
The selected dog and lighthouse examples illustrate coherent subjects and
scene structure, rather than typical-case performance. The teddy-bear row
also exposes limited color variation despite plausible texture.
Other saved samples expose limitations: the proposal omits the person in a
dining-room caption, does not reliably separate requested teddy-bear colors,
and renders a requested upside-down stop sign upright. Distribution-level
FID and CLIP improvements therefore do not establish reliable counting,
spatial relations or complete prompt adherence. The teacher reference also
uses a much larger inference budget.

\section{Limitations and conclusion}
\paragraph{Resource and training-budget constraints.}
Available GPU resources and project budget limit each student to 36,000
generator updates, effective batch 16, rank-32 adapters and one training seed.
This amounts to 576,000 samples in the generator training stream. Published
training budgets vary: SwiftBrush reports approximately 65,000 updates at
batch 64~\cite{nguyen2024swiftbrush}, while the original consistency-model
experiments use 600,000--1,000,000 iterations across different datasets and
training variants~\cite{song2023consistency}. These settings differ in model
capacity, batch size and training task, so raw iteration counts do not equalize
compute. Our constrained runs may leave both the proposal and baselines
under-optimized; longer training could change their relative performance.
The results establish a comparison within our budget, rather than the
performance ceiling of the published methods.

The completed results support stronger full-validation quality for the
combined proposal than our DMD baseline at a matched total update
budget. The factorial study favors perceptual regression alone for FID;
the added diversity term does not establish a consistent benefit.
One training seed provides no estimate
of run-to-run uncertainty. Adapter capacity, teacher-pretraining overlap,
the restricted regression bank and validation reuse during development limit
generality. Perceptual distance can increase through artifacts and cannot
guarantee semantic diversity. Broader claims require additional seeds and controlled
compute measurements.

{\small
\bibliographystyle{ieeenat_fullname}
\bibliography{references}
}
\end{document}
